\documentclass[10pt,a4paper]{amsart}
\usepackage[margin=19mm]{geometry}
\usepackage{amsmath,amssymb,mathtools}
\usepackage[scr=rsfs,cal=euler]{mathalfa}
\usepackage{xfrac}
\usepackage[unicode,hidelinks,bookmarks=false]{hyperref}
\newcommand{\ie}{i.e.\ }
\newcommand{\cf}{cf.\ }
\newcommand{\resp}{respectively\ }
\newcommand{\Subsection}{\subsection}
\providecommand{\nobreakdash}{\nobreak\hbox{-}}
\setlength{\emergencystretch}{2em}
\multlinegap0pt
\usepackage{bm}
\usepackage{bbm}
\usepackage{dsfont}
\usepackage{xspace}
\usepackage{cancel}
\usepackage{mathtools}
\usepackage{amsthm}
\makeatletter
\@ifundefined{theorem}{\newtheorem{theorem}{Theorem}}{}
\@ifundefined{lemma}{\newtheorem{lemma}[theorem]{Lemma}}{}
\makeatother
\title[An improved bound for the strong clique index of graphs]{An improved bound for the strong clique index of graphs}
\author{Hitesh Kumar, Bojan Mohar, Shivaramakrishna Pragada}
\date{Source: arXiv:2607.02698v1 (CC BY 4.0)}
\begin{document}
\maketitle

\noindent\emph{Source and licence.} An improved bound for the strong clique index of graphs. Version arXiv:2607.02698v1 (CC BY 4.0), distributed under the Creative Commons Attribution 4.0 International licence, as recorded in the arXiv licence metadata for this exact version and shown on the abstract page https://arxiv.org/abs/2607.02698v1. Full rights notice: licenses/arxiv-2607.02698-CC-BY-4.0.txt.\par\smallskip

\noindent\textit{Editorial scope.} Counted unit: the Lemma whose source label is \texttt{lemma:bipartite\_subgraph} in the article (source0.tex lines 252--258, source label \texttt{lemma:bipartite\_subgraph}), delivered together with the article's own complete original proof of it (source0.tex lines 260--382, one \texttt{proof} environment of 123 lines). No mathematics is written, repaired or re-derived: statement and proof are the article's own text line for line. Required context retained uncounted: lemma:beta (lines 216--218). Every definition, display and label of those lines is retained unchanged. Omitted from this package and not redistributed: 1--215, 219--251, 259--259, 383--610. Nothing inside a retained excerpt is omitted or altered: every retained body is delivered exactly as the article's lines are written, so no omission receipt of any kind accompanies this package. Citations: the retained excerpts carry no citation command at all, so no bibliography entry of the article is transcribed and no reference is redistributed. Reference adaptation: none was needed and none was made; every reference inside the delivered unit resolves inside it, so no unresolved reference or citation is produced. Printed slips kept literal: no printed word, symbol, hypothesis or number of the counted statement or of its proof was altered. Layout and class: the article's own class is \texttt{article} with packages including fullpage, amsmath, amsthm, amsfonts, amssymb, amstext, mathrsfs, enumerate, graphicx, ragged2e, lscape, framed, xcolor, subfiles; the wrapper is the lane's pinned \texttt{amsart} editorial wrapper at 10pt on A4, which restates the theorem-like environments the retained text needs and adds the article's own macro definitions that the retained text needs (none), with extra wrapper packages (bm, bbm, dsfont, xspace, cancel, mathtools). The delivered numbering is the unit's own. Assets: no retained span contains a figure environment, a graphics-inclusion command, a PDF-page inclusion, a table or a TikZ picture; no image file is copied into this package, the package declares no asset dependency and no image file is redistributed with it. Licence: version arXiv:2607.02698v1 of this article is distributed under the Creative Commons Attribution 4.0 International licence, as recorded in the arXiv licence metadata for this exact version and shown on the abstract page https://arxiv.org/abs/2607.02698v1. A full rights notice is delivered at licenses/arxiv-2607.02698-CC-BY-4.0.txt. Characters: every retained excerpt is pure ASCII, so the delivered engine, XeLaTeX with its default Latin Modern text font, reports no missing character.
\begin{lemma}\label{lemma:beta} We have 
    \[ \frac{1}{4} < \beta^* \le \frac{620}{1987} < 0.3120282 < \frac{5}{16}.\]
\end{lemma}

\begin{lemma}\label{lemma:bipartite_subgraph}
Let $G$ be a graph and let $H$ be a bipartite subgraph of $G$. If
$E(H)$ is a clique in $L(G)^2$, then
\[
   |E(H)|\le \beta^*\,\sigma_G(H)^2.
\]
\end{lemma}

\begin{proof}
Consider the subgraph $G[V(H)]$ of $G$ induced by the vertices of $H$. It is clear that $E(H)$ is also a clique in $L(G[V(H)])^2$ and $\sigma_{G[V(H)]}(H) \le \sigma_G(H)$. Thus, we can replace $G$ with $G[V(H)]$; equivalently, we can assume that $H$ is a spanning subgraph of $G$. Throughout, let
\[
   \sigma:=\sigma_G(H)\quad \text{and} \quad \Delta:=\Delta(H).
\]
We proceed by induction on $|E(H)|$. If $|E(H)|=1$, then $\sigma\ge 2$, which implies
\[
   |E(H)|=1\le 4\beta^*\le \beta^*\,\sigma^2,
\]
since $\beta^* > \frac{1}{4}$. 

So, assume $|E(H)|\ge 2$. Choose $v\in V(H)$ such that $\deg_H(v)=\Delta$. Let $X\cup Y$ be a bipartition
of $H$ and assume $v\in X$. Define
\[
   A:=N_H(v)\subseteq Y,  \quad  C:= N_G(v)\setminus A, \quad B:=V(H)\setminus N_G[v], \quad  F:=H[B].
\]
Clearly, every edge in $E(H)\setminus E(F)$ is incident with a vertex in $N_G[v]$. Therefore
\[
   |E(H)\setminus E(F)|\le |E_H(A)| + |E_H(C)|,
\]
where $E_H(S)$ denotes the set of $H$-edges incident with $S\subseteq V(H)$.
Therefore, 
\begin{equation}\label{eq:H_bound_1}
   |E(H)|\le |E_H(A)| + |E_H(C)| + |E(F)|. 
\end{equation}
In what follows, we will estimate $|E_H(A)|$ and $|E_H(C)|$ using structural arguments and $|E(F)|$ by induction hypothesis, thus obtaining the desired bound for $|E(H)|$.

First,
\begin{equation}\label{eq:C_bound}
   |E_H(C)|\le \Delta |C| =\Delta(\deg_G(v)-\Delta).
\end{equation}

Next, we estimate $|E_H(A)|$. Note that for each $a\in A$, the edge $va$ lies in $H$, so $\deg_G(v)+\deg_G(a)\le \sigma.$
Thus
\[
   \sum_{a\in A}\deg_G(a)\le \Delta(\sigma-\deg_G(v)).
\]
Define 
\[
  B_X:=B\cap X,\quad B_Y:=B\cap Y, \quad   \Pi:= |E_G(A,B_Y)| + |E_G(A,B_X)\setminus E(H)|,
\]
where $E_G(S,R)$ denotes the edges of $G$ with one endpoint in $S$ and the other in $R$, where $S, R\subseteq V(G)$. The term $\Pi$ counts the $G$-edges incident with $A$ and $B$, which consume $G$-degree at vertices of $A$ but are not counted as $H$-edges incident with $A$. Thus, 
\begin{equation}\label{eq:A_bound}
    |E_H(A)| \le \sum_{a\in A}\deg_G(a)-\Pi \le \Delta(\sigma-\deg_G(v))-\Pi.
\end{equation}
Note that 
\[
   \Pi= \sum_{y\in B_Y} |N_G(y)\cap A| + \sum_{x\in B_X} |(N_G(x)\cap A)\setminus N_H(x)|.
\]
Take an edge $xy\in E(F)$, where $x\in B_X$ and $y\in B_Y$. For every
$a\in A$, the two $H$-edges $va$ and $xy$ are at distance at most
two in $L(G)$ since $E(H)$ forms a clique in $L(G)^2$. Since $x,y\notin N_G[v]$, this can only happen through
an edge $ax$ or an edge $ay$. Therefore,
\[
   A\subseteq N_G(x)\cup N_G(y) \implies   |A|\le |N_G(x)\cap A| + |N_G(y)\cap A|.
\]
Since $|A| = \Delta$, we get 
\[ \Delta \le  |N_G(y)\cap A|  + |(N_G(x)\cap A)\setminus N_H(x)| + |N_H(x)\cap A|. \]
For $x\in B$, define 
\[ h_x : = |N_H(x)\cap A|, \quad f_x: = \deg_{F}(x).\]
Summing over all $xy\in E(F)$ and using the fact that $f_x, f_y\le \Delta$, we get
\begin{align*}
    \Delta |E(F)|
    & \le  \sum_{y\in B_Y}|N_G(y)\cap A|f_y  + \sum_{x\in B_X}|(N_G(x)\cap A)\setminus N_H(x)|f_x + \sum_{x\in B_X}h_xf_x\\
   & \le  \left(\sum_{y\in B_Y}|N_G(y)\cap A|  + \sum_{x\in B_X}|(N_G(x)\cap A)\setminus N_H(x)|\right)\Delta + \sum_{x\in B_X}h_xf_x\\
   &  = \Pi \Delta +  \sum_{x\in B_X}h_xf_x
\end{align*}
for any $xy\in E(F)$, which implies
\begin{equation}\label{eq:Pi_bound}
   |E(F)|\le \Pi+\frac{1}{\Delta}\sum_{x\in B_X}h_xf_x. 
\end{equation}
Combining \eqref{eq:H_bound_1}, \eqref{eq:C_bound}, \eqref{eq:A_bound} and \eqref{eq:Pi_bound}, we get 
\begin{equation}\label{eq:H_bound_2}
   |E(H)| \le \Delta(\sigma-\Delta) + \frac{1}{\Delta}\sum_{x\in B_X}h_xf_x. 
\end{equation}
Note that 
\begin{equation}\label{eq:h_x_estimate}
  \sum_{x\in B_X} h_x  = |E_H(A, B_X)| \le \sum_{a\in A} \deg_H(a)\le |A|\Delta = \Delta^2.  
\end{equation}
Also, $F$ is a bipartite subgraph of $G[B]$ such that $E(F)$ induces a clique in $L(G[B])^2$. Moreover, as discussed before, for any edge $xy\in E(F)$, every vertex in $A$ is adjacent to either $x$ or $y$. This means
\[ \deg_{G[B]}(x) + \deg_{G[B]}(y) \le \deg_G(x) + \deg_{G}(y) - \Delta \le \sigma - \Delta,\]
which implies 
\[ \sigma_{G[B]}(F) \le \sigma - \Delta.\]
Since $|E(F)| < |E(H)|$, by induction hypothesis,
\begin{equation}\label{eq:f_x_estimate}
   \sum_{x\in B_X}f_x = |E(F)|\le \beta^*\ \sigma_{G[B]}(F)^2 \le \beta^* (\sigma - \Delta)^2. 
\end{equation}
Define 
\[ p: = \frac{\beta^*(\sigma-\Delta)^2}{\Delta^2 + \beta^*(\sigma-\Delta)^2} \quad \text{and} \quad q: = 1-p = \frac{\Delta^2}{\Delta^2 +\beta^*(\sigma-\Delta)^2}.\]
For any $x\in B_X$, we have $h_x + f_x \le \Delta$, which implies 
\[
   \frac{h_xf_x}{\Delta}
   \le
   \frac{h_xf_x}{h_x+f_x},
\]
with the convention that the right side is $0$ whenever $h_x + f_x = 0$. Since $p + q = 1$, it is easy to check that 
\[    \frac{h_xf_x}{h_x+f_x} \le p^2 h_x + q^2 f_x.\]
Indeed, 
\[ (p^2 h_x + q^2 f_x)(h_x + f_x) - h_xf_x = (ph_x - qf_x)^2 \ge 0.\]
Summing over $x\in B_X$, we get
\[ \frac{1}{\Delta}\sum_{x\in B_X}h_xf_x \le  p^2 \Delta^2 + q^2 \beta^* (\sigma - \Delta)^2 = \frac{\Delta^2\beta^*(\sigma-\Delta)^2}{\Delta^2 + \beta^*(\sigma-\Delta)^2},\]
by \eqref{eq:h_x_estimate} and \eqref{eq:f_x_estimate}. Hence, by \eqref{eq:H_bound_2}, we have
\[
   |E(H)| \le \Delta(\sigma-\Delta) + \frac{\Delta^2\beta^*(\sigma-\Delta)^2}{\Delta^2 + \beta^*(\sigma-\Delta)^2}.
\]
Set $t=\Delta/\sigma$. Then
\[
   \frac{|E(H)|}{\sigma^2}
   \le
   t(1-t)+
   \frac{\beta^* t^2(1-t)^2}{t^2+\beta^*(1-t)^2}.
\]
Observe that since $|E(H)|>0$, we have $t\in (0,1)$. Now, by Lemma \ref{lemma:beta}, 
\[P_{\beta^*}\left(\frac{t}{1-t}\right)\ge 0\] for all $t<1$. After simplification, one can see that this inequality is equivalent to 
\[t(1-t)+
   \frac{\beta^* t^2(1-t)^2}{t^2+\beta^*(1-t)^2}\le \beta^*. 
\]
Therefore,
\[
   |E(H)|\le \beta^*\sigma^2,
\]
as required.
\end{proof}

\end{document}
